% =============================================================================
% Chapter 7: Fault Tolerance and Reliability — ML Systems Lecture Slides
% =============================================================================
% Volume II, Chapter 7. Prerequisites: Fleet Stack, C3, compute infra,
% network fabrics, storage, distributed training, collective communication.
%
% IMAGES (SVG -> PDF):
%   - failure-spectrum.pdf        - reliability-gap.pdf
%   - young-daly-tradeoff.pdf     - checkpoint-timeline.pdf
%   - elastic-training.pdf        - training-vs-serving-ft.pdf
%   - graceful-degradation.pdf    - circuit-breaker-states.pdf
%   - observability-pillars.pdf
% =============================================================================
\documentclass[aspectratio=169, 12pt]{beamer}
\usepackage{../../assets/beamerthememlsys}

\mlsyssetup{
  volume       = {Volume II},
  chapter      = {Chapter 7},
  logo         = {../../assets/img/logo-mlsysbook.png},
  instlogo     = {../../assets/img/logo-harvard.png},
  chaptertitle = {Fault Tolerance and Reliability},
}

% --- Fonts ---
\usepackage{fontspec}
\setsansfont{Helvetica Neue}[
  BoldFont={Helvetica Neue Bold},
  ItalicFont={Helvetica Neue Italic},
  BoldItalicFont={Helvetica Neue Bold Italic},
]
\setmonofont{JetBrains Mono}[Scale=0.85]

% --- Packages ---
\usepackage{booktabs}
\usepackage{amsmath}

% --- Image paths ---
\graphicspath{
  {images/}
}

% --- Chapter-specific macros ---
\newcommand{\MTBF}{\text{MTBF}}

% --- Helper: safe image include ---
\newcommand{\safeimg}[2][width=\textwidth,keepaspectratio]{%
  \IfFileExists{images/#2}{\includegraphics[#1]{#2}}{%
    \IfFileExists{#2}{\includegraphics[#1]{#2}}{%
      \fbox{\parbox[c][2.5cm][c]{0.85\linewidth}{\centering\footnotesize\textcolor{midgray}{[Missing image]}}}%
    }%
  }%
}

% --- Section count for navigation ---
\setcounter{mlsystotalsections}{8}

\title{Fault Tolerance and Reliability}
\author{Vijay Janapa Reddi}
\institute{Harvard University}
\date{}

\begin{document}

% =============================================================================
% TITLE SLIDE
% =============================================================================
\mlsystitle{Fault Tolerance}{The Immune System of Infrastructure}{cover_fault_tolerance.png}

% =============================================================================
% LEARNING OBJECTIVES
% =============================================================================
\begin{frame}{Learning Objectives}
\note{[2 min] Walk through objectives. Emphasize that this chapter transitions
from ``how to coordinate GPUs'' (Ch 5--6) to ``how to survive when GPUs die.''
Ask: ``How many of you have lost training progress to a hardware crash?''}

\small
\begin{enumerate}
  \item Calculate system-level \textbf{MTBF} from component failure rates and derive optimal checkpoint intervals using the \textbf{Young-Daly formula}
  \item Design distributed \textbf{checkpoint and recovery} strategies for multi-terabyte model states across thousands of GPUs
  \item Implement serving \textbf{redundancy}, \textbf{replication}, and \textbf{failover} within millisecond latency requirements
  \item Apply \textbf{graceful degradation}: model fallback, feature fallback, and load shedding
  \item Construct \textbf{observability infrastructure} using metrics, logs, and traces to diagnose non-deterministic failures
  \item Evaluate fault tolerance trade-offs between \textbf{training} and \textbf{serving} workloads
\end{enumerate}

\end{frame}

\begin{frame}{Visual Language}
\note{[1 min] Explain the semantic color system used throughout the course.
These colors are consistent across all diagrams and slides.}

\small
Throughout this course, colors carry meaning:

\vspace{0.3cm}
\begin{columns}[T]
  \begin{column}{0.45\textwidth}
    \begin{mlsyscard}{computestroke}
    \textbf{Blue} --- Compute / Processing\\
    {\footnotesize GPU ops, forward/backward pass, inference}
    \end{mlsyscard}
    \vspace{0.15cm}
    \begin{mlsyscard}{datastroke}
    \textbf{Green} --- Data / Memory\\
    {\footnotesize Data flow, caches, healthy paths}
    \end{mlsyscard}
  \end{column}
  \begin{column}{0.45\textwidth}
    \begin{mlsyscard}{routingstroke}
    \textbf{Orange} --- Routing / Scheduling\\
    {\footnotesize Load balancers, batch windows}
    \end{mlsyscard}
    \vspace{0.15cm}
    \begin{mlsyscard}{errorstroke}
    \textbf{Red} --- Error / Cost / Bottleneck\\
    {\footnotesize Loss, decode phase, waste}
    \end{mlsyscard}
  \end{column}
\end{columns}
\end{frame}



% =============================================================================
\section{Failure at Scale}
% =============================================================================

\begin{frame}{The ML Failure Spectrum}
\note{[3 min] Three failure categories, each requiring different detection and
recovery. Key insight: silent corruption is the most dangerous because no error
is raised. Ask: ``Which type is hardest to debug?''
Common error: students think all failures are hard crashes.}

% --- Layout: FULL-WIDTH diagram ---
\centering
\safeimg[width=0.92\textwidth,height=4.8cm,keepaspectratio]{failure-spectrum.pdf}

\vspace{0.15cm}
\mlsysinsight{Key:}{Hard failures are easy. Silent data corruption is the real threat at scale.}

\end{frame}

\begin{frame}{The Mathematics of Inevitable Failure}
\note{[3 min] Walk through the reliability product rule. The key equation:
system MTBF = component MTBF / N. This is the central insight of the chapter.
Ask: ``If one GPU fails once per year, how often does a 10K cluster fail?''
Common error: students think adding redundancy always helps.}

\small
\begin{columns}[T]
  \begin{column}{0.52\textwidth}
    Individual component reliability:
    $$R_{single}(t) = e^{-\lambda t}$$

    \vspace{0.1cm}
    For $N$ identical components (series):
    $$R_{system}(t) = e^{-N\lambda t}$$

    \vspace{0.1cm}
    \begin{mlsyscard}{crimson}
    \textbf{System MTBF:}
    $$\MTBF_{system} = \frac{\MTBF_{component}}{N}$$
    \end{mlsyscard}
  \end{column}
  \begin{column}{0.44\textwidth}
    \vspace{0.2cm}
    \mlsysconcept{Scale inverts reliability:}{}

    \vspace{0.1cm}
    {\footnotesize
    \renewcommand{\arraystretch}{1.15}
    \begin{tabular}{@{}rr@{}}
      \toprule
      \textbf{GPUs} & \textbf{Cluster MTBF} \\
      \midrule
      8       & 260 days \\
      1,000   & 2 days \\
      10,000  & \textbf{5 hours} \\
      25,000  & \textbf{2 hours} \\
      \bottomrule
    \end{tabular}
    }

    \vspace{0.1cm}
    {\footnotesize\textcolor{midgray}{Per-GPU MTBF: 50,000 hours}}
  \end{column}
\end{columns}

\end{frame}

\begin{frame}{The Reliability Gap}
\note{[2 min] This figure shows cluster MTBF dropping below training duration
as clusters grow. The ``danger zone'' is where training cannot complete without
fault tolerance. Meta's Llama 3: 419 failures in 54 days at 16K GPUs.
If short: focus on the 10K GPU data point.}

% --- Layout: FULL-WIDTH image ---
\centering
\safeimg[width=0.92\textwidth,height=4.8cm,keepaspectratio]{reliability-gap.pdf}

\vspace{0.1cm}
\mlsysalert{Implication:}{At 10K+ GPUs, fault tolerance is not optional --- it is physically necessary.}

\end{frame}

% --- ACTIVE LEARNING: Micro-Retrieval Cue ---
\begin{frame}{Quick Check}
\note{[1 min] Answer: 50000/10000 = 5 hours.}

\centering
\vspace{1.0cm}
{\Large\bfseries Quick Check}

\vspace{0.6cm}
{\large At 10,000 GPUs with 50,000h per-GPU MTBF,\\what is the cluster MTBF?}

\vspace{0.5cm}
{\normalsize\textcolor{midgray}{15 seconds --- then I will cold-call.}}

\end{frame}



% --- ACTIVE LEARNING 1: Predict ---
\begin{frame}{Predict: The 9s of Reliability}
\note{[2 min] Give students 60 seconds. Answer: $(0.99999)^{10000} \approx 0.90$.
Only 90\% chance the entire cluster is up for one hour! This primes the
checkpoint discussion. Ask 2--3 students to share answers.}

\centering
\vspace{0.8cm}
{\Large\bfseries Think--Write--Share}

\vspace{0.5cm}
{\large Each GPU has \textbf{99.999\%} uptime.\\
You have \textbf{10,000 GPUs}.\\[0.3cm]
What is the probability the \emph{entire cluster}\\
survives \textbf{one hour} without a failure?}

\vspace{0.5cm}
{\normalsize Write your estimate. \textcolor{midgray}{(60 seconds)}}

\end{frame}

% =============================================================================
\section{The Young-Daly Law}
% =============================================================================

\mlsysfocus{The Young-Daly Law}{%
$\tau_{\text{opt}} = \sqrt{2 \cdot T_{\text{write}} \cdot \MTBF}$%
}

\begin{frame}{Optimal Checkpoint Interval}
\note{[3 min] Walk through the U-shaped cost curve. Left side: too frequent
checkpoints waste I/O time. Right side: too infrequent checkpoints waste
rework after failure. The minimum is the Young-Daly optimum.
Ask: ``Which side is worse?''}

% --- Layout: FULL-WIDTH diagram ---
\centering
\safeimg[width=0.88\textwidth,height=4.0cm,keepaspectratio]{young-daly-tradeoff.pdf}

\vspace{0.1cm}
\small
\begin{columns}[T]
  \begin{column}{0.30\textwidth}
    \centering
    \textcolor{computestroke}{\textbf{Too frequent}}\\
    {\scriptsize I/O overhead dominates}
  \end{column}
  \begin{column}{0.30\textwidth}
    \centering
    \textcolor{datastroke}{\textbf{Optimal: $\tau_{opt}$}}\\
    {\scriptsize Minimizes total waste}
  \end{column}
  \begin{column}{0.30\textwidth}
    \centering
    \textcolor{errorstroke}{\textbf{Too rare}}\\
    {\scriptsize Rework cost dominates}
  \end{column}
\end{columns}

\end{frame}

\begin{frame}{Young-Daly: Worked Example}
\note{[3 min] Walk through step by step. Key: the formula gives ~27 minutes,
much more precise than intuition. Note the square-root relationship: halving
MTBF only increases frequency by 1.4x. If short: just show final result.}

\small
\textbf{Scenario: Llama-3 training on 16,000 H100 GPUs}

\vspace{0.15cm}
\renewcommand{\arraystretch}{1.15}
{\footnotesize
\begin{tabular}{@{}lll@{}}
  \toprule
  \textbf{Parameter} & \textbf{Value} & \textbf{Source} \\
  \midrule
  Checkpoint write time ($T_w$) & 2 min  & Storage BW \\
  Cluster MTBF                  & 3 hrs (180 min)  & 16K GPUs \\
  \bottomrule
\end{tabular}
}

\vspace{0.15cm}
$$\tau_{\text{opt}} = \sqrt{2 \cdot 2 \cdot 180} = \sqrt{720} \approx \mathbf{27\text{ min}}$$

\vspace{0.15cm}
\begin{mlsyscard}{datastroke}
{\footnotesize
\textbf{Checkpoint every 10 min}: 17\% overhead (too frequent)\\
\textbf{Checkpoint every 60 min}: 15\% overhead (too risky)\\
\textbf{Checkpoint every 27 min}: \textbf{7\% overhead} (optimal)}
\end{mlsyscard}

\end{frame}

% =============================================================================
\section{Checkpointing}
% =============================================================================

\begin{frame}{Checkpoint-Recovery Timeline}
\note{[3 min] Walk through the Gantt chart. Key phases: productive training,
checkpoint writes, failure event, lost work, recovery. The ``lost work'' is
on average tau/2. Total failure cost = lost work + recovery latency.
Ask: ``What determines the recovery latency?''}

% --- Layout: FULL-WIDTH diagram ---
\centering
\safeimg[width=0.92\textwidth,height=4.5cm,keepaspectratio]{checkpoint-timeline.pdf}

\vspace{0.1cm}
\mlsysalert{Total failure cost:}{Lost work ($\approx\tau/2$) + Recovery latency (load + warmup). Restart can take 3--5$\times$ checkpoint time.}

\end{frame}

\begin{frame}{Checkpoint State: What Gets Saved?}
\note{[2 min] A full GPT-3 checkpoint is 2.1 TB. That is the bottleneck.
Adam optimizer state (m + v) doubles the naive size. The checkpoint storm
problem: 10K GPUs writing simultaneously. If short: focus on the size.}

\footnotesize
\begin{columns}[T]
  \begin{column}{0.50\textwidth}
    \textbf{Full training state:}
    \begin{itemize}\setlength\itemsep{0pt}
      \item Model weights (FP16): 2 bytes/param
      \item Adam $m$ (first moment): 4 bytes/param
      \item Adam $v$ (second moment): 4 bytes/param
      \item Data loader position, LR schedule
      \item Random seeds for reproducibility
    \end{itemize}

    \vspace{0.1cm}
    \textbf{Total}: 12 bytes/param

    \vspace{0.1cm}
    \begin{mlsyscard}{errorstroke}
    {\footnotesize \textbf{GPT-3} (175B params):\\
    Full checkpoint $\approx$ \textbf{2.1 TB}}
    \end{mlsyscard}
  \end{column}
  \begin{column}{0.46\textwidth}
    \textbf{The Checkpoint Storm:}

    \vspace{0.1cm}
    When 10,000 GPUs write simultaneously:
    \begin{itemize}\setlength\itemsep{0pt}
      \item Network fabric saturated
      \item Parallel file system overwhelmed
      \item I/O contention cascades
    \end{itemize}

    \vspace{0.1cm}
    \textbf{Mitigation strategies:}
    \begin{itemize}\setlength\itemsep{0pt}
      \item \textcolor{computestroke}{Async checkpointing}: write in background
      \item \textcolor{datastroke}{Tiered staging}: GPU $\to$ CPU $\to$ NVMe $\to$ PFS
      \item \textcolor{routingstroke}{Sharded writes}: each rank saves its shard
    \end{itemize}
  \end{column}
\end{columns}

\end{frame}

% --- ACTIVE LEARNING 2: Exercise ---
\begin{frame}{Your Turn: Checkpoint Interval}
\note{[3 min] Give students 90 seconds. Answer: tau = sqrt(2 * 5 * 240) =
sqrt(2400) = 49 min. Note: intuitive ``every hour'' is close but suboptimal.
``Every 15 min'' (common guess) wastes 2x more than optimal.}

\small
\begin{columns}[T]
  \begin{column}{0.55\textwidth}
    {\normalsize\bfseries Calculate the Optimal Interval}

    \vspace{0.2cm}
    Your 1,000-GPU cluster has:
    \begin{itemize}
      \item $T_{\text{write}}$ = 5 min (checkpoint save)
      \item $\MTBF$ = 4 hours (240 min)
    \end{itemize}

    \vspace{0.15cm}
    \textbf{Apply} the Young-Daly formula to find $\tau_{\text{opt}}$.

    \vspace{0.1cm}
    {\footnotesize Compare your answer to ``every hour'' and ``every 15 min.''}

    \vspace{0.1cm}
    {\footnotesize\textcolor{midgray}{(90 seconds --- then compare with a neighbor)}}
  \end{column}
  \begin{column}{0.42\textwidth}
    \pause
    \begin{mlsyscard}{datastroke}
    \textbf{Solution:}\\[0.1cm]
    {\footnotesize
    $\tau = \sqrt{2 \times 5 \times 240}$\\
    $\tau = \sqrt{2400} \approx$ \textbf{49 min}\\[0.1cm]
    Every 60 min: $\approx$13\% overhead\\
    Every 15 min: $\approx$36\% overhead\\
    Every 49 min: \textcolor{datastroke}{\textbf{8\% overhead}}
    }
    \end{mlsyscard}
  \end{column}
\end{columns}

\end{frame}

% =============================================================================
\section{Elastic Training}
% =============================================================================

\begin{frame}{Elastic Training: Dynamic Recovery}
\note{[3 min] Static training aborts on failure; elastic training adapts.
The key: convert hard failures into throughput degradation. This requires
batch size adjustment, LR scaling, and state resharding.
Ask: ``Why not just wait for a replacement GPU?''}

% --- Layout: FULL-WIDTH diagram ---
\centering
\safeimg[width=0.88\textwidth,height=4.5cm,keepaspectratio]{elastic-training.pdf}

\vspace{0.1cm}
\mlsysinsight{Key:}{Elastic training converts fatal failures into recoverable throughput drops.}

\end{frame}

\begin{frame}{Elastic Training: Mechanisms}
\note{[2 min] Four mechanisms make elasticity work. LR scaling is critical:
linear rule (Goyal 2017) or square-root rule. Gradient accumulation maintains
effective batch size. Spot instances enable 50\%+ cost savings.
If short: focus on batch size and LR.}

\footnotesize
\begin{columns}[T]
  \begin{column}{0.48\textwidth}
    \textbf{Adapting to $N-1$ workers:}

    \vspace{0.1cm}
    \textbf{1. Batch size adjustment}\\
    {\scriptsize Reduce global batch or increase per-worker load}

    \vspace{0.1cm}
    \textbf{2. Learning rate scaling}
    $$\eta_{new} = \eta_{base} \times \sqrt{\frac{N_{new}}{N_{base}}}$$
    {\scriptsize Square-root rule for conservative adjustment}

    \vspace{0.1cm}
    \textbf{3. Gradient accumulation}\\
    {\scriptsize Double micro-batches if workers halved}

    \vspace{0.1cm}
    \textbf{4. State resharding}\\
    {\scriptsize Redistribute ZeRO/FSDP shards}
  \end{column}
  \begin{column}{0.48\textwidth}
    \textbf{Spot instance arbitrage:}

    \vspace{0.1cm}
    \begin{mlsyscard}{routingstroke}
    {\footnotesize Cloud spot instances: 60--90\% discount\\
    \textbf{Caveat}: Can be reclaimed anytime\\[0.1cm]
    Elastic training transforms preemption from fatal error to recoverable resize}
    \end{mlsyscard}

    \vspace{0.1cm}
    {\footnotesize
    \renewcommand{\arraystretch}{1.1}
    \begin{tabular}{@{}lll@{}}
      \toprule
      \textbf{Framework} & \textbf{Elastic} & \textbf{Reshard} \\
      \midrule
      TorchElastic & Yes & Manual \\
      DeepSpeed    & Yes & Auto \\
      Ray Train    & Yes & Auto \\
      Horovod      & Yes & Manual \\
      \bottomrule
    \end{tabular}
    }
  \end{column}
\end{columns}

\end{frame}

% =============================================================================
\section{Serving Resilience}
% =============================================================================

\begin{frame}{Training vs Serving Fault Tolerance}
\note{[3 min] The most important comparison in this chapter. Training tolerates
minutes of recovery; serving demands milliseconds. Training state is terabytes;
serving state is gigabytes (KV cache). Same principle, different latency budgets.
Ask: ``Which is harder to make fault-tolerant?''}

% --- Layout: FULL-WIDTH diagram ---
\centering
\safeimg[width=0.92\textwidth,height=4.8cm,keepaspectratio]{training-vs-serving-ft.pdf}

\vspace{0.1cm}
{\small Same principle: redundancy + detection + recovery. Different latency budget: \textcolor{computestroke}{\textbf{minutes}} vs \textcolor{routingstroke}{\textbf{milliseconds}}.}

\end{frame}

\begin{frame}{Stateless vs Stateful Serving}
\note{[2 min] Stateless serving is easy: retry on another replica. Stateful
serving (LLM KV cache) is hard: state loss forces reprocessing all previous
turns. The distinction determines entire fault tolerance architecture.}

\footnotesize
\renewcommand{\arraystretch}{1.15}
\begin{tabular}{@{}lp{4cm}p{4.5cm}@{}}
  \toprule
  \textbf{Aspect} & \textbf{Stateless} & \textbf{Stateful} \\
  \midrule
  \textbf{Examples}   & Image classification, embedding generation & LLM chat (KV cache), streaming ASR \\
  \textbf{Routing}    & Any replica & Session-affine replica \\
  \textbf{On failure} & Retry on another replica & Potential state loss \\
  \textbf{Recovery}   & Restart + load weights (seconds) & Reload + reconstruct context (seconds--min) \\
  \textbf{Redundancy} & Active-active replicas & Replicated state + standby \\
  \bottomrule
\end{tabular}

\vspace{0.15cm}
\begin{mlsyscard}{routingstroke}
{\footnotesize \textbf{Availability with replication}: $A_{system} = 1 - (1 - A_{single})^R$.
Two replicas at 99\% each $\to$ 99.99\%. Three replicas $\to$ 99.9999\%.}
\end{mlsyscard}

\end{frame}

% =============================================================================
\section{Graceful Degradation}
% =============================================================================

\begin{frame}{Graceful Degradation: Model Fallback}
\note{[3 min] When primary systems are overwhelmed, do not crash --- degrade.
The hierarchy: primary GPU model > secondary CPU model > cached static results.
War story: e-commerce fallback lost only 2\% revenue for 3 hours.
Ask: ``When is 75\% quality acceptable?''}

% --- Layout: FULL-WIDTH diagram ---
\centering
\safeimg[width=0.88\textwidth,height=4.8cm,keepaspectratio]{graceful-degradation.pdf}

\vspace{0.1cm}
{\small \textcolor{errorstroke}{\textbf{75\%}} quality at \textcolor{datastroke}{\textbf{100\%}} availability beats \textcolor{datastroke}{\textbf{100\%}} quality at \textcolor{errorstroke}{\textbf{0\%}} availability.}

\end{frame}

\begin{frame}{Circuit Breaker Pattern}
\note{[2 min] Three states: closed (normal), open (fail fast), half-open (probe).
The key protection: prevents one failing service from exhausting the entire
fleet's connection pool. Named after electrical circuit breakers.
If short: focus on the open state preventing cascades.}

\small
\begin{columns}[T]
  \begin{column}{0.48\textwidth}
    \safeimg[width=\textwidth,height=5.5cm,keepaspectratio]{circuit-breaker-states.pdf}
  \end{column}
  \begin{column}{0.48\textwidth}
    \textbf{Three states:}

    \vspace{0.1cm}
    \textcolor{datastroke}{\textbf{Closed}} --- Normal flow\\
    {\footnotesize Requests pass to dependency}

    \vspace{0.1cm}
    \textcolor{errorstroke}{\textbf{Open}} --- Fail fast\\
    {\footnotesize Reject immediately, prevent cascade}

    \vspace{0.1cm}
    \textcolor{routingstroke}{\textbf{Half-Open}} --- Probe\\
    {\footnotesize Test if dependency recovered}

    \vspace{0.15cm}
    \begin{mlsyscard}{errorstroke}
    {\footnotesize \textbf{Without circuit breakers}: one failing model replica exhausts all connection pools $\to$ cascade failure across entire fleet.}
    \end{mlsyscard}
  \end{column}
\end{columns}

\end{frame}

% --- ACTIVE LEARNING 3: Discussion ---
\begin{frame}{Discussion: Which Failure Costs More?}
\note{[3 min] Turn-and-talk. Students discuss in pairs for 90 seconds.
Cold-call 2--3 pairs. Point: training failures waste compute (\$\$\$), but
serving failures lose users and revenue in real time. No single right answer.
If short: do a show-of-hands poll.}

\centering
\vspace{0.6cm}
{\large\bfseries Turn and Talk \textcolor{midgray}{(90 seconds)}}

\vspace{0.4cm}
{\normalsize A 10,000-GPU cluster crashes during training.\\
Meanwhile, the serving fleet drops to 50\% capacity for 10 minutes.\\[0.2cm]
\alert{Which incident costs more?}}

\vspace{0.3cm}
\begin{columns}[c]
  \begin{column}{0.40\textwidth}\centering\small
    \textcolor{computestroke}{\textbf{Training crash}}\\
    Hours of lost compute
  \end{column}
  \begin{column}{0.40\textwidth}\centering\small
    \textcolor{routingstroke}{\textbf{Serving degradation}}\\
    Lost user trust + revenue
  \end{column}
\end{columns}

\end{frame}

% =============================================================================
\section{Observability}
% =============================================================================

\begin{frame}{Three Pillars of Observability}
\note{[3 min] Metrics tell WHEN, Logs tell WHAT, Traces tell WHY. All three
must be correlated to diagnose distributed failures. Heisenbugs: adding
instrumentation changes timing, masking the race condition.
Ask: ``Why can't you just add print statements?''}

% --- Layout: FULL-WIDTH diagram ---
\centering
\safeimg[width=0.88\textwidth,height=4.5cm,keepaspectratio]{observability-pillars.pdf}

\vspace{0.1cm}
\mlsysconcept{Heisenbugs:}{Adding logging changes timing --- the bug disappears when you try to observe it.}

\end{frame}

\begin{frame}{Training Failure Signatures}
\note{[2 min] Three patterns to recognize: loss spike (transient), divergence
(permanent corruption), hang (deadlock). Each requires a different response.
Loss spikes may be hardware bit flips, not bad data batches.}

\footnotesize
\begin{columns}[T]
  \begin{column}{0.48\textwidth}
    \textbf{Pattern recognition for training:}

    \vspace{0.15cm}
    \textcolor{routingstroke}{\textbf{Loss Spike + Recovery}}\\
    {\scriptsize Transient data issue or bit flip\\
    Response: Monitor, investigate if repeated}

    \vspace{0.1cm}
    \textcolor{errorstroke}{\textbf{Loss Divergence}}\\
    {\scriptsize Silent corruption or desync\\
    Response: \textbf{Rollback} to earlier checkpoint}

    \vspace{0.1cm}
    \textcolor{computestroke}{\textbf{Loss Hang (Flatline)}}\\
    {\scriptsize Deadlock in collective comm.\\
    Response: Timeout detection, restart}
  \end{column}
  \begin{column}{0.48\textwidth}
    \textbf{Serving failure signatures:}

    \vspace{0.15cm}
    \textcolor{routingstroke}{\textbf{Latency Spikes}}\\
    {\scriptsize Resource contention, cold cache\\
    Response: Scale up, warm cache}

    \vspace{0.1cm}
    \textcolor{errorstroke}{\textbf{Error Rate Increase}}\\
    {\scriptsize Dependency failure, format change\\
    Response: Circuit breaker, investigate}

    \vspace{0.1cm}
    \textcolor{computestroke}{\textbf{Silent Quality Decay}}\\
    {\scriptsize Model drift, feature degradation\\
    Response: Distribution monitoring}

    \vspace{0.1cm}
    \textcolor{errorstroke}{\textbf{Cascade Failure}}\\
    {\scriptsize One failure propagates via retries\\
    Response: Circuit breakers + isolation}
  \end{column}
\end{columns}

\end{frame}

% =============================================================================
\section{Case Studies}
% =============================================================================

\begin{frame}{Meta OPT-175B: Training Through Failure}
\note{[3 min] 992 A100 GPUs, 2 months, 35+ hardware failures. Async
checkpointing reduced overhead from 12\% to 3\%. ``Last good checkpoint''
rollback for loss spikes. Maintained >50\% MFU despite daily interruptions.
If short: focus on async checkpointing win.}

\footnotesize
\begin{columns}[T]
  \begin{column}{0.50\textwidth}
    \textbf{OPT-175B Training (Meta, 2022):}
    \begin{itemize}\setlength\itemsep{0pt}
      \item 992 A100 GPUs, 2 months
      \item 35+ significant hardware failures
      \item System MTBF $<$ 24 hours
    \end{itemize}

    \vspace{0.1cm}
    \textbf{Key innovations:}
    \begin{itemize}\setlength\itemsep{0pt}
      \item 5-min heartbeat timeout detection
      \item Async checkpointing: 12\% $\to$ \textbf{3\%} overhead
      \item ``Last good checkpoint'' rollback for loss spikes
    \end{itemize}
  \end{column}
  \begin{column}{0.46\textwidth}
    \begin{mlsyscard}{computestroke}
    {\footnotesize \textbf{Result:} Sustained $>$50\% MFU despite daily interruptions.}
    \end{mlsyscard}

    \vspace{0.1cm}
    \begin{mlsyscard}{errorstroke}
    {\footnotesize \textbf{Netflix Chaos Engineering:}\\
    Injected failures in ML serving. Discovered 12 unknown failure modes. MTTR: 45 min $\to$ \textbf{8 min}.\\[0.05cm]
    Popularity fallback maintained 94\% of recommendation quality.}
    \end{mlsyscard}
  \end{column}
\end{columns}
\mlsyscite{Zhang et al., ``OPT,'' 2022; Meta Llama 3 Report, 2024}

\end{frame}

\begin{frame}{Google TPU Pods \& DeepSpeed Elasticity}
\note{[2 min] Two contrasting approaches. Google: replace entire ``slice''
(hardware unit of failure). DeepSpeed: redistribute shards across remaining
workers. Both achieve <3\% overhead. Key lesson: the atomic unit of failure
determines the recovery strategy.}

\footnotesize
\begin{columns}[T]
  \begin{column}{0.48\textwidth}
    \textbf{Google TPU v4 Pods:}
    \begin{itemize}\setlength\itemsep{0pt}
      \item 4,096 chips in 3D toroidal mesh
      \item One chip failure $\to$ topology hole
      \item Strategy: \textbf{Slice Swap}
    \end{itemize}

    \vspace{0.1cm}
    {\scriptsize Don't repair running topology --- replace the entire slice. ``Hot pool'' of standby capacity. Recovery target: $<$3 minutes. Maintains $>$98\% effective throughput.}

    \vspace{0.1cm}
    \mlsysconcept{Lesson:}{Hardware grouping is the atomic failure unit.}
  \end{column}
  \begin{column}{0.48\textwidth}
    \textbf{Microsoft DeepSpeed:}
    \begin{itemize}\setlength\itemsep{0pt}
      \item ZeRO-3: state sharded across GPUs
      \item Node failure = lost 1/N of state
      \item Strategy: \textbf{Elastic resizing}
    \end{itemize}

    \vspace{0.1cm}
    {\scriptsize Detect via heartbeat ($<$2\% overhead). Redistribute state to $N-1$ workers. Adjust gradient accumulation for batch consistency. Enables training on spot instances.}

    \vspace{0.1cm}
    \mlsysconcept{Lesson:}{Software masks hardware volatility.}
  \end{column}
\end{columns}

\end{frame}

% =============================================================================
\section{Wrap-Up}
% =============================================================================

\begin{frame}{Fallacies}
\note{[2 min] Three quantitative rebuttals. Hardware failures are only 15--25\%
of incidents. Silent corruption is real: 0.01--0.1\% of GPUs per year.
Elastic training needs checkpoints too.}

\footnotesize
\textbf{Fallacy:} \textit{Hardware failures are the main concern.}\\
Software bugs cause 30--40\% of incidents, config errors 20--30\%. Hardware is only 15--25\%. Hardware redundancy alone leaves most failure modes unaddressed.

\vspace{0.1cm}
\textbf{Fallacy:} \textit{Silent data corruption is rare in modern hardware.}\\
0.01--0.1\% of GPUs per year exhibit silent corruption not caught by ECC. For 10K GPUs: 1--10 silent events/year, hundreds of bit flips monthly. Gradient NaNs blamed on LR may be bit flips.

\vspace{0.1cm}
\textbf{Fallacy:} \textit{Elastic training eliminates the need for checkpointing.}\\
Elasticity cannot survive below-minimum worker count, accumulated degradation, or software bugs in remaining workers. Elastic + checkpoint is the answer.

\end{frame}

\begin{frame}{Pitfalls}
\note{[2 min] Setting intervals by intuition can be 2--3x off.
Ignoring restart overhead: 5-min save + 20-min restart = 25-min cost.
Testing fault tolerance only during real failures is too late.}

\footnotesize
\textbf{Pitfall:} \textit{Setting checkpoint interval by intuition.}\\
{\scriptsize ``Every hour'' may be 2--3$\times$ off. For MTBF=4h, $T_w$=5min: $\tau_{opt}=49$ min. If $T_w$=15 min: $\tau_{opt}=85$ min, not ``every 15 min.''}

\vspace{0.1cm}
\textbf{Pitfall:} \textit{Ignoring restart overhead in checkpoint planning.}\\
{\scriptsize 5-min save + 20-min restart = 25-min effective cost, not 5. Use $T_{save} + T_{restart}$ in the formula.}

\vspace{0.1cm}
\textbf{Pitfall:} \textit{Testing fault tolerance only during real failures.}\\
{\scriptsize Checkpoint restoration and circuit breakers may be untested. \textbf{Chaos engineering} transforms ``we think it works'' into ``we know it works.''}

\end{frame}

% --- RETRIEVAL PRACTICE ---

% --- MUDDIEST POINT ---
\begin{frame}{Muddiest Point}
\note{[2 min] Quick anonymous poll. Students write on a slip of paper or submit
digitally. Collect and scan for patterns. Address the top 2--3 confusions in the
next lecture's opening. This closes the feedback loop.}

\centering
\vspace{1.0cm}
{\Large\bfseries What was the \alert{muddiest point} today?}

\vspace{0.8cm}
{\normalsize Write down the concept you found \textbf{most confusing}.}

\vspace{0.5cm}
{\small\textcolor{midgray}{Anonymous. One sentence. Submit before you leave.}}

\end{frame}

\begin{frame}{What Were the Key Ideas?}
\note{[2 min] Retrieval practice. Students write 90 seconds, no notes.
Do NOT show next slide yet. Walk around the room.}

\centering
\vspace{1.5cm}
{\Large\bfseries Close your notes.}

\vspace{0.8cm}
{\large Write down the \textbf{4 most important concepts} from today.}

\vspace{0.8cm}
{\normalsize\textcolor{midgray}{90 seconds --- no peeking.}}

\end{frame}

\begin{frame}{Key Takeaways}
\note{[2 min] Reveal. Walk through each bullet. Emphasize: MTBF/N is the
central equation, Young-Daly governs checkpointing, elasticity converts
failures to throughput drops, training vs serving need different strategies.}

\scriptsize
\begin{itemize}\setlength\itemsep{0pt}
  \item \textbf{Scale guarantees failure}: $\MTBF_{sys} = \MTBF_{comp}/N$. A 10K-GPU cluster fails every 5 hours.
  \item \textbf{Young-Daly governs checkpoints}: $\tau_{opt} = \sqrt{2 \cdot T_w \cdot \MTBF}$. Typical: 20--30 min.
  \item \textbf{Checkpoint storm}: GPT-3 = 2.1 TB. Async + tiered staging prevents saturation.
  \item \textbf{Elastic training}: Convert fatal failures to throughput drops via batch/LR/shard adjustment.
  \item \textbf{Training vs serving}: Minutes of recovery vs millisecond failover.
  \item \textbf{Graceful degradation}: 75\% quality at 100\% availability $>$ 100\% quality at 0\%.
  \item \textbf{Observability}: Metrics (when) + Logs (what) + Traces (why) = distributed debugging.
\end{itemize}

\end{frame}

\begin{frame}{References}
\note{[1 min] Point students to canonical papers.}

\small
\mlsysref{Young74}{J.W. Young. ``A First-Order Checkpoint Interval.'' CACM, 1974.}
\mlsysref{Daly06}{J. Daly. ``A Higher Order Estimate of Checkpoint Intervals.'' FGCS, 2006.}
\mlsysref{Zhang+22}{Zhang et al. ``OPT: Open Pre-trained Transformers.'' arXiv, 2022.}
\mlsysref{Kokolis+25}{Kokolis et al. ``Characterizing GPU Failures in the Wild.'' HPCA, 2025.}
\mlsysref{Meta24}{Meta. ``The Llama 3 Herd of Models.'' arXiv, 2024.}
\mlsysref{Rasley+20}{Rasley et al. ``DeepSpeed: System Optimizations.'' KDD, 2020.}

\end{frame}

\begin{frame}{Next Lecture: Fleet Orchestration}
\note{[1 min] Forward hook. We can survive failures now --- but who decides
which jobs run where? Scheduling, preemption, multi-tenancy, capacity
planning. The ``brain'' of the fleet.}

\footnotesize
\begin{columns}[c]
  \begin{column}{0.30\textwidth}
    \centering
    \begin{mlsyscard}{computestroke}
    {\large\bfseries Scheduling}\\
    {\footnotesize Slurm, K8s,\\gang scheduling}
    \end{mlsyscard}
  \end{column}
  \begin{column}{0.30\textwidth}
    \centering
    \begin{mlsyscard}{routingstroke}
    {\large\bfseries Multi-Tenancy}\\
    {\footnotesize Fair sharing,\\preemption, quotas}
    \end{mlsyscard}
  \end{column}
  \begin{column}{0.30\textwidth}
    \centering
    \begin{mlsyscard}{datastroke}
    {\large\bfseries Capacity}\\
    {\footnotesize Planning, bin-packing,\\resource allocation}
    \end{mlsyscard}
  \end{column}
\end{columns}

\vspace{0.2cm}
\centering
{\normalsize A fault-tolerant cluster still needs a \textbf{brain}.}\\[0.1cm]
{\small Who runs what, where, and when?}

\end{frame}



\appendix

\begin{frame}{Backup: Young-Daly Sensitivity Analysis}
\note{Backup: How checkpoint interval changes with parameters.}

\footnotesize
\textbf{$\tau_{\text{opt}} = \sqrt{2 \cdot T_w \cdot \text{MTBF}}$}

\vspace{0.15cm}
\renewcommand{\arraystretch}{1.15}
\begin{tabular}{@{}rrrrr@{}}
  \toprule
  \textbf{$T_w$ (min)} & \textbf{MTBF=1h} & \textbf{MTBF=4h} & \textbf{MTBF=12h} & \textbf{MTBF=24h} \\
  \midrule
  1  & 11 min & 22 min & 38 min & 54 min \\
  5  & 24 min & 49 min & 85 min & 120 min \\
  15 & 42 min & 85 min & 147 min & 208 min \\
  \bottomrule
\end{tabular}

\vspace{0.1cm}
{\footnotesize Halving $T_w$ reduces optimal interval by $\sqrt{2} \approx 1.4\times$, not $2\times$.}

\end{frame}

\begin{frame}{Backup: Async Checkpointing Deep Dive}
\note{Backup: How async checkpointing works in practice.}

\footnotesize
\textbf{Three-phase async checkpoint:}
\begin{enumerate}\setlength\itemsep{1pt}
  \item \textbf{Snapshot}: Copy GPU state to pinned CPU memory (ms)
  \item \textbf{Background write}: CPU writes to NVMe/PFS (seconds)
  \item \textbf{Replication}: Async copy to remote storage (minutes)
\end{enumerate}

\vspace{0.1cm}
\textbf{Training pause = Phase 1 only} ($<$100 ms for most models).

\vspace{0.1cm}
\begin{mlsyscard}{datastroke}
{\footnotesize OPT-175B: Async reduced checkpoint overhead from \textbf{12\% to 3\%} of training time.}
\end{mlsyscard}

\end{frame}


\end{document}
